% Conclusion — Securing LLMs Against Prompt Injection
% Academic manuscript section (IEEE / arXiv style)

\section{Conclusion}
\label{sec:conclusion}

This paper studied model-agnostic pre-inference sanitization for direct prompt injection.
We compared an unprotected baseline and two lexical filters with a hybrid handcrafted context-aware gate (Method~A) and a learned soft-fusion variant (Method~B), under a two-layer protocol that separates Bypass Rate from True ASR and reports false positives, robustness, ablations, and latency.

Lexical defenses alone proved insufficient on the held-out test partition once end-to-end compliance was measured: they reduced Bypass Rate only modestly and left True ASR unchanged relative to the unprotected baseline.
Both context-aware methods improved security substantially.
Method~A achieved the strongest protection, including zero successful attacks on the static test set, but incurred higher false-positive rates, especially on technical edge-case prompts.
Method~B offered a milder usability cost while leaving residual True ASR.
Under adaptive stylistic rewriting, all active defenses degraded, yet the same ordering was preserved.

Ablation analysis clarifies where these gains originate.
Method~A should be understood as a hybrid defense rather than a pure ensemble scorer: semantic hard triggers provide the majority of protection, while weighted multi-signal scoring serves as a complementary layer for borderline cases.
Method~B's gains depend primarily on semantic features; removing them returns performance to unprotected baseline levels on this partition.
We therefore retain the present architecture and report it transparently, rather than redesigning the sanitizer around an inaccurate ensemble narrative.

Within the stated scope---direct textual injection, a compact curated corpus, and a single open instruction-tuned target---the results support three practical conclusions.
First, portable context-aware sanitization can outperform standard regex and keyword controls when security is measured end-to-end.
Second, hybrid hard-trigger designs can deliver stronger static protection than soft fusion alone, at an explicit usability cost.
Third, semantic evidence is central to both handcrafted and learned variants in this setting.

Future work may reduce Method~A's edge-case over-blocking, explore adaptive or learned fusion that balances signal contributions more evenly, strengthen adaptive attackers, and extend evaluation to indirect and agentic benchmarks.
Even under the present constraints, the study provides an evidence-oriented account of a deployable pre-inference sanitizer and of the architectural choices that drive its defensive capability.
